\documentclass[11pt]{article}
\usepackage[margin=1in]{geometry}
\usepackage{amsmath,amssymb,amsthm}
\usepackage{booktabs,enumitem,xcolor,array,tabularx}
\usepackage{microtype}
\usepackage[colorlinks=true,linkcolor=blue!50!black,citecolor=blue!50!black,urlcolor=blue!50!black]{hyperref}
\usepackage[round]{natbib}
\newcolumntype{Y}{>{\raggedright\arraybackslash}X}
\setlength{\emergencystretch}{2em}
\newtheorem{theorem}{Theorem}
\newcommand{\E}{\mathbb{E}}
\newcommand{\pending}[1]{\textcolor{red!60!black}{[pending: #1]}}
\hypersetup{pdftitle={Which Rollouts Are Worth Generating?}}
\title{\textbf{Which Rollouts Are Worth Generating?}\\[4pt]
\large Exact Bias, Variance, and Cost of Adaptive Rollout Collection in RLVR}
\author{Mohamed Ayman\\\small Working draft, October 16, 2026. Results marked \pending{} are running.}
\date{}
\begin{document}
\maketitle

\begin{abstract}
Reinforcement learning with verifiable rewards spends most of its compute generating answers, and practical systems decide adaptively how many answers to generate per prompt: they drop groups whose answers all agree, abandon prompts after a short probe, or sample until a correct answer appears. Each rule changes what training optimizes, how noisy each update is, and what it costs, yet rules are compared through heuristics and proxies. We show that for binary rewards, the first two moments of the update produced by any rule that decides from the rewards seen so far, including rules with random group size, depend on only four numbers per prompt: its success probability, the squared norm of the gradient of that probability, and the mean squared gradient norms of correct and of wrong answers. This makes the bias, variance, and cost of every such rule exactly computable, and makes rollout collection an optimization problem that we solve exactly. Under this accounting, DAPO's dynamic sampling has exactly the per-answer efficiency of the fixed group it filters, sample-until-correct rules are the least efficient rules in use, and the optimal rule's advantage over the best rule in use depends on how gradient signal varies with difficulty. Measuring the four numbers in Qwen2.5-Math-1.5B shows that the signal fades for very easy and very hard prompts, which makes the optimal rule robust and puts its predicted gain over the best rule in use at 1.6 to 2.2 times, and that wrong answers carry about 2.5 times larger squared gradients than correct ones, which the theory turns into a one-line change to GRPO: scale failure advantages by the measured ratio. In a pre-registered pilot that change did not improve accuracy. We show why a better rule can be worth nothing: without the orthogonality assumption, the progress of any rule at a budget of $B$ answers per step is governed by a critical budget, the rule's gradient noise scale. Below it a rule with $k$ times the efficiency is worth $k$ times the answers (about $\sqrt k$ under Adam at a fixed learning rate); above it, no reweighting can help. All the quantities involved can be measured inside a training run. \pending{budget ladder and in-training measurements}
\end{abstract}

\section{Introduction}
Generating answers dominates the cost of RL with verifiable rewards, so a family of rules now decides adaptively how many answers each prompt receives. DAPO drops groups whose answers all agree and refills the batch \citep{yu2025dapo}; SARA abandons prompts after a short probe \citep{sara2026}; Reinforce-Ada and related methods sample until enough correct answers appear \citep{xiong2025reinforceada}; allocation methods move a budget toward prompts that look informative \citep{knapsack2025,yao2025gvm}. Each such rule does three things at once: it changes the implicit objective, it changes the update's noise, and it changes the cost. Existing comparisons optimize one of these through a proxy.

For binary rewards, recent work shows that fixed-size, reward-only weightings optimize an explicit function of each prompt's success probability \citep{davis2025objective,mroueh2025grpo,thrampoulidis2025shaping}, and that dropping all-agreeing groups rescales the update without changing its direction \citep{okonkwo2026arcus}. These results concern the mean of the update for fixed group sizes. The rules practitioners use have random group sizes, and the quantity that decides their value is the update's noise per unit cost.

\textbf{Contributions.}
\begin{enumerate}[leftmargin=*,itemsep=2pt,topsep=2pt]
\item \textbf{A two-moment theorem} (Section~\ref{sec:theory}). For any rule that stops and assigns advantages from the rewards observed so far, the mean update is $c(p)\nabla p$ and the second moment is $F_1D_1(p)+F_0D_0(p)+\|\nabla p\|^2X(p)$, where $F_1,F_0$ are the mean squared gradient norms of correct and wrong answers. Random group sizes need no special treatment.
\item \textbf{Exact accounting and the optimal rule} (Section~\ref{sec:optimal}). An exact lattice engine evaluates any rule, and dynamic programming finds the rule that maximizes training value per generated answer. We compare the rules in use against it.
\item \textbf{Measurements in a real model} (Section~\ref{sec:measure}). Wrong answers carry about 2.5 to 2.8 times larger squared gradients than correct ones in Qwen2.5-Math-1.5B, both in the last layer and over all parameters, while being only 9\% longer.
\item \textbf{The critical budget} (Section~\ref{sec:regime}). Dropping orthogonality, the exact one-step progress of a rule depends on its gradient noise scale in answers; better rules pay off only below it, and the payoff is measurable during training.
\item \textbf{Tests} (Section~\ref{sec:training}). A pre-registered pilot of scaled failure advantages (no accuracy gain), and a pre-registered budget ladder with in-training measurements of the critical budget and of each rule's efficiency. \pending{result}
\end{enumerate}

\section{The two-moment theorem}\label{sec:theory}
A prompt has policy $\pi$, score $\ell(y)=\nabla\log\pi(y)$, binary reward $R(y)$, and success probability $p$. A reward-adapted rule draws answers independently, stops at a bounded stopping time $N$ of the reward sequence, and returns $g=\sum_{i\le N}A_i\ell(y_i)$ with each $A_i$ a function of $(r_1,\dots,r_N)$; filtering and abandonment are $A=0$. Let $\kappa(1)=1/p$, $\kappa(0)=-1/(1-p)$, and $F_r=\E[\|\ell\|^2\mid R=r]$.

\begin{theorem}
For every reward-adapted rule, $\E g=c(p)\nabla p$ with $c(p)=\E\sum_{i\le N}A_i\kappa(r_i)$, and
\[
\E\|g\|^2=F_1D_1(p)+F_0D_0(p)+\|\nabla p\|^2X(p),
\]
with $D_r=\E\sum_{i\le N:\,r_i=r}A_i^2$ and $X=\E\sum_{i\neq j\le N}A_iA_j\kappa(r_i)\kappa(r_j)$. The functions $c,D_0,D_1,X$ and the expected cost $\E N$ depend only on the rule and $p$.
\end{theorem}
\emph{Proof.} $\E[\ell\,1\{R=1\}]=\nabla p$ and $\E\ell=0$ give $\E[\ell\mid R=r]=\kappa(r)\nabla p$. Couple the rule to an infinite i.i.d.\ sequence; given the reward sequence, $N$ and the $A_i$ are constants and the answers are independent with that conditional mean and conditional second moment $F_{r_i}$. Take expectations over the rewards, whose law depends only on $p$. $\square$

The structure breaks under token-level length normalization, where weights depend on lengths rather than rewards, under non-binary rewards, and under clipping or stale samples; exact counterexamples are in the code. All statements are checked by exact enumeration over softmax policies, including two rules with random group size.

\section{Exact accounting and the optimal rule}\label{sec:optimal}
\textbf{Criterion.} Prompts come from a difficulty distribution and gradients of different prompts are treated as orthogonal. With the step size chosen optimally, the improvement in mean accuracy per generated answer is proportional to $\mathrm{Eff}=\mathrm{Sig}^2/(\mathrm{Cost}\cdot\mathrm{M2})$, where $\mathrm{Sig}=\E[c(p)\|\nabla p\|^2]$, $\mathrm{M2}$ is the mean second moment, and $\mathrm{Cost}$ the mean number of answers. Because rescaling advantages leaves $\mathrm{Eff}$ unchanged, the optimal rule is found by dynamic programming over the lattice of (wrong, correct) counts after a Dinkelbach transform; a solved policy, re-evaluated by the independent lattice engine, matches to six digits.

\textbf{Rules in use.} Table~\ref{tab:catalog} gives efficiency relative to the optimal rule in the main model, $\|\nabla p\|^2\propto p^2(1-p)^2$ with equal $F_0$ and $F_1$ in the high-noise regime that real models occupy. Three conclusions are exact within the model. DAPO's dynamic sampling has the per-answer efficiency of the fixed group it filters, because dropped groups cost their answers and contribute nothing. Sample-until-correct rules are the least efficient, spending budget on hard prompts that carry little gradient. A probe that abandons all-agreeing prompts and fills the rest is the best rule in use.

\begin{table}[h]
\centering\small
\begin{tabularx}{\linewidth}{@{}>{\raggedright\arraybackslash}p{0.42\linewidth}YYYY@{}}
\toprule
\textbf{Rule} & \textbf{Uniform} & \textbf{Mostly hard} & \textbf{U-shaped} & \textbf{Mostly easy}\\
\midrule
Fixed group of 8 (GRPO, Dr.\ GRPO, DAPO) & 0.61 & 0.38 & 0.40 & 0.38\\
Fixed group of 16 & 0.65 & 0.41 & 0.43 & 0.41\\
Fixed group of 8, std-normalized & 0.57 & 0.35 & 0.36 & 0.35\\
Probe 4, abandon if all agree, fill to 16 & 0.77 & 0.60 & 0.67 & 0.60\\
Sample until 1 correct, cap 16 & 0.26 & 0.09 & 0.12 & 0.30\\
Sample until 2 correct and 2 wrong, cap 16 & 0.41 & 0.22 & 0.22 & 0.22\\
\bottomrule
\end{tabularx}
\caption{Efficiency relative to the optimal rule, by difficulty distribution of the prompts.}
\label{tab:catalog}
\end{table}

\textbf{Structure of the optimum.} The optimal rule samples at least three answers, continues while the minority outcome stays above about a quarter of the answers, and stops when outcomes are lopsided; it assigns advantages state by state, nearly zeroing lopsided low-signal groups. A two-parameter band rule that copies only the stopping region reaches just 18\% to 59\% of the optimum, so the advantage sizing matters as much as the stopping.

\textbf{Sensitivity.} The size of the gain depends on how $\|\nabla p\|^2$ varies with difficulty: the optimal rule beats the best rule in use by 1.29 to 1.66 times when the signal fades as $p^2(1-p)^2$, by 1.15 to 1.33 times for $p(1-p)$, and by 2.3 to 4.2 times if it is flat. A policy solved under either fading shape keeps at least 96\% of the optimum under the other, but falls to 10\% to 31\% if the truth is flat, and no mixture policy is robust to both. The shape must therefore be measured before the optimal rule is deployed (Section~\ref{sec:measure}).

\section{When is a better rule worth anything? The critical budget}\label{sec:regime}
The criterion above treats prompt gradients as orthogonal, so it counts progress only on the prompts in the batch. Held-out progress comes from the part of each prompt's gradient that is shared with the population gradient $\nabla J=\E_x\nabla p_x$, and then the cross-prompt terms matter. Let a step use $P$ prompts drawn independently, each processed by the same rule, with update $g=\frac1P\sum_x u_x$. Write $\mu=\E u=\E_x[c(p_x)\nabla p_x]$, $s=\|\mu\|^2$, $m_2=\E\|u\|^2$ (given prompt by prompt by Theorem~1), $V=m_2-s$, $n=\E N$, and $B=Pn$ for the answers per step.

\begin{theorem}
$\E g=\mu$ and $\E\|g\|^2=s+nV/B$. For $J$ that is $L$-smooth along the step, the best SGD step gives one-step progress
\[
\Delta(B)=\frac{\langle\nabla J,\mu\rangle^2}{2L\,(s+nV/B)}.
\]
\end{theorem}
\emph{Proof.} Independence across prompts; then maximize $\eta\langle\nabla J,\mu\rangle-\frac{L}{2}\eta^2\E\|g\|^2$ over $\eta$. $\square$

The \emph{critical budget} $B^\star=nV/s$, the rule's gradient noise scale in answers \citep{mccandlish2018empirical}, separates two regimes. Below it, progress is linear in $B$ and proportional to the per-answer efficiency $\langle\nabla J,\mu\rangle^2/(n\,m_2)$, which reduces to the criterion of Section~\ref{sec:optimal} when prompt gradients are orthogonal; a rule with $k$ times the efficiency is worth $k$ times the answers. Above it, $\Delta\to\|\nabla J\|^2\cos^2(\nabla J,\mu)/(2L)$: progress depends neither on the budget nor on the rule's noise, only on the direction of its mean update, and reweighting prompts can only lose, because the unweighted estimator points exactly along $\nabla J$. The ratio of two rules at any budget,
\[
\frac{\Delta_R(B)}{\Delta_0(B)}=\frac{\langle\nabla J,\mu_R\rangle^2}{\langle\nabla J,\mu_0\rangle^2}\cdot\frac{s_0+n_0V_0/B}{s_R+n_RV_R/B},
\]
involves only quantities that can be measured during training from two independent halves of a batch and an independent unbiased reference for $\nabla J$ (Section~\ref{sec:training}).

\textbf{Adam.} With a fixed learning rate and Adam's second-moment estimate dominated by noise, the update is roughly the gradient divided by its per-coordinate noise level. Cutting the noise variance $k$-fold then raises first-order progress per step by about $\sqrt k$, while the step size and hence the second-order penalty stay fixed. A rule with $k$ times the efficiency is then worth about $\sqrt k$ times the progress per step at a fixed learning rate, and the full $k$ only with a larger learning rate, as in the square-root scaling rule for adaptive methods \citep{malladi2022sde}. For the scaled-failure rule this turns the predicted 1.12 to 1.56 times into 1.06 to 1.25 times, a small effect at a fixed learning rate.

\section{Measuring the four numbers}\label{sec:measure}
\textbf{Method.} In Qwen2.5-Math-1.5B \citep{yang2024qwen25math}, we sample answers to MATH training problems on-policy and compute exact per-answer gradients in the output projection and the last-layer MLP, using the single-backward identity of the companion STOP study. $F_1$ and $F_0$ are averages of squared answer-gradient norms; $\|\nabla p\|^2$ is estimated without bias from pairs, since $\E[(r_i-c)(r_j-c)\langle\ell_i,\ell_j\rangle]=\|\nabla p\|^2$ for independent answers and any $c$ independent of the pair.

\textbf{Gradient asymmetry.} Over 40 prompts with 16 answers each, $F_1=4.2\times10^5$ (90\% interval $[3.6,4.8]\times10^5$) and $F_0=1.17\times10^6$ ($[1.04,1.30]\times10^6$), so $F_1/F_0=0.36$, while wrong answers are only 9\% longer (421 versus 386 tokens). Over all 1.5B parameters, calibration before training gives $F_1/F_0=0.395$ from 29 correct and 99 wrong answers. Wrong answers carry larger gradients because their tokens are less likely.

\textbf{Signal against difficulty.} $\|\nabla p\|^2$ is small against per-answer noise (ratio of order $10^3$ to $10^4$), and 16 answers per prompt do not resolve its shape; heavy tails from a few long wrong answers make naive intervals too narrow. With 96 answers per prompt (4{,}560 pairs) in the MLP block, over 24 prompts spread across difficulty, the signal is resolved. On prompts with success rate at most 0.15 or at least 0.85 the pooled estimate is $0.73\pm0.60$; on prompts between 0.3 and 0.7 it is $15.5\pm5.4$. The ratio, 0.05, rejects a flat shape (predicted 1) and disfavors $p(1-p)$ (0.36); a weighted fit prefers $p^2(1-p)^2$ (chi-square per prompt 0.35, against 0.54 for $p(1-p)$ and 0.77 for flat). The signal fades for very easy and very hard prompts, which is the regime where the optimal rule is robust (Section~\ref{sec:optimal}). With the measured gradient ratio, the best rule in use reaches 0.46 to 0.63 of the optimum, so the optimal rule's predicted gain is 1.6 to 2.2 times.

\section{Training tests}\label{sec:training}
\textbf{Prediction.} When $F_0>F_1$, the optimal advantages weight a failure at about $F_1/F_0$ relative to a success on mixed groups, whatever the shape of $\|\nabla p\|^2$. Applied to a fixed group of 8 with the measured ratio, this alone raises efficiency by 1.12 to 1.56 times in 11 of 12 combinations of shape and difficulty distribution; the exception is mostly easy prompts with flat signal (0.85). It is a one-line change to GRPO.

\textbf{Design (pre-registered).} Qwen2.5-Math-1.5B base, 3{,}920 numeric-answer MATH training problems, on-policy updates with one gradient step per batch, 16 prompts by 8 answers, 150 steps, AdamW at $10^{-6}$. The arms differ only in advantages, Dr.\ GRPO's $r-\bar r$ \citep{liu2025drgrpo} versus the same with failures scaled by $w=F_1/F_0$ measured before training, so both generate the same number of answers. Two seeds, each running both arms on one L40S. Primary endpoint: greedy accuracy on the 364 numeric MATH-500 problems after 150 steps.

\textbf{Interim results.} Calibration gives $w=0.395$ and $0.420$ in the two seeds. The baseline arm finished at 47.5\% (seed 1) and 46.7\% (seed 2) greedy accuracy, up from 13.2\%. In seed 1 the scaled arm finished at 45.3\%, 2.2 points below the baseline, after being tied at steps 50 and 100 (41.5\% versus 41.8\%, 42.9\% versus 42.0\%). Its training reward was higher in each of the first four 25-step windows (0.296 versus 0.287, 0.484 versus 0.465, 0.550 versus 0.540, 0.571 versus 0.563) and then tied (0.594 versus 0.592, 0.577 versus 0.581); the mean reward over the last 50 steps is 0.586 in both arms, and answer lengths match (407 tokens). Both arms use the same random numbers at each step, so these are paired comparisons. A single-seed difference of 2.2 points on 364 problems is within evaluation noise (the standard error of one arm's accuracy on 364 problems is 2.6 points), so the seed-1 result is consistent with no effect on final accuracy and a small early speedup. Seed 2 shows the same pattern so far (41.2\% versus 40.9\% at step 50, 42.9\% versus 44.2\% at step 100). \pending{seed 2 final; pre-registered verdict}

\textbf{Reading the pilot through Section~\ref{sec:regime}.} A null result at a fixed learning rate is what the theory predicts if training at 128 answers per step is near or above the critical budget, or if Adam reduces the gain to its square root. The design below measures both instead of assuming them.

\textbf{Budget ladder and in-training diagnostics (pre-registered).} Per seed, Dr.\ GRPO runs at 32, 64, 128, and 256 answers per step, the optimal rule at 64 and 128 (blocks of 4, a stop-or-continue decision from the solved lookup table after each block, cap 16, the table's per-state advantages, re-solved every 10 steps from a beta-binomial fit to the last 200 prompts, prompts per step chosen so that cumulative answers track the budget), and the scaled-failure rule at 128. At every step, prompts are split at random into halves $A$ and $B$; with $g_A,g_B$ the arm's gradients from each half and $r_B$ the Dr.\ GRPO estimator on the first 4 answers of each prompt in $B$, an unbiased direction for $\nabla J$ independent of $g_A$, we record $\|g_A\|^2$, $\|g_B\|^2$, $\langle g_A,g_B\rangle$, and $\langle g_A,r_B\rangle$. These give $s$, $V$, $B^\star$, and $\langle\nabla J,\mu\rangle$ up to a common scale, so the progress ratio of Theorem~2 is computed from each run's own measurements and compared with its observed accuracy curve. Accuracy is measured greedily on the 364 MATH-500 problems at ten points, sampled at temperature 1 (the quantity the theory optimizes) at three points, and greedily on 3{,}382 further MATH test problems at the end, all saved per problem for paired comparisons. \pending{14 runs started October 17}

\section{Related work}
The implicit objectives of fixed-size reward-only weightings are given by \citet{davis2025objective}, with GRPO, Dr.\ GRPO, pass@k, and maximum-likelihood variants worked out in \citet{mroueh2025grpo,thrampoulidis2025shaping,zheng2026rl2ml,tajwar2026maxrl,heckel2026asym}. ARCUS proves the score-conditioning identity and the dynamic-sampling corollary \citep{okonkwo2026arcus}. RL2ML gives the covariance of success-only estimators with fixed group size \citep{zheng2026rl2ml}. Allocation and stopping methods optimize proxies or bounds \citep{yao2025gvm,xiong2025reinforceada,sara2026,knapsack2025}. Our results add random group sizes, the full second moment with both $F_0$ and $F_1$, the exact optimal rule, and measurements of the four numbers in a real model.

\section{Limitations}
The efficiency criterion is one step of progress on mean accuracy; Theorem~2 removes the orthogonality assumption but keeps the one-step, smooth, SGD view, and the Adam statement is a noise-dominated approximation. Exploration effects are outside the theory. Cost is counted in answers, not tokens. The shape of $\|\nabla p\|^2$ is measured in one parameter block of one model, and the training comparison is a two-seed pilot.

\bibliographystyle{plainnat}
\bibliography{refs}
\end{document}
